\section{Discussion}\label{sec:discussion}
\paragraph{Detection and repair are separate problems (exploratory).} Apart from flagging every step, read-back is the
only trigger that flags silent write faults reliably, and it does so on all three models; the
uncertainty trigger catches a few by chance. Whether a flag turns
into success depends on the model. A dropped write is repaired by repeating the call, which Haiku
usually does and Qwen3-8B sometimes does, while Qwen3-4B almost never gets there. A wrong-value or
partial write needs the agent to find what changed and undo it, which only Haiku does often. This
matches reports that, under repeated poisoning, agents adopt wrong answers even after
checking~\cite{anon2026tools}, and that a monitor which issues corrective directives improves
recovery~\cite{kaynar2026tools}. Telling the agent \emph{what} differed from the expected state may
help the small models, as telling a client the exact fault kind does elsewhere~\cite{woo2026same}; we did not test this, to keep the NOTICE identical
across triggers.

\paragraph{More checks can hurt (exploratory).} On Haiku, adding the contract to read-back lowers
success on wrong and partial writes (\RHaikuWrongBothSuccess{} and \RHaikuPartialBothSuccess{}, against
\RHaikuWrongReadbackSuccess{} and \RHaikuPartialReadbackSuccess{} with read-back alone). These cells
hold \RHaikuWrongReadbackN{} and \RHaikuPartialReadbackN{} episodes, and every pair where only one of
the two arms succeeded favours read-back alone. An inspection of those pairs suggests a cause. Both arms
raise the same first flag, on the faulted write. With both checks, the agent's later re-reads of the
corrupted record fail the contract's type check (wrong writes) or schema check (partial writes) and
are flagged as well, and in the episodes we read the agent then ends the task without
repairing the record. Our hypothesis is that the NOTICE, which says the \emph{result} may not match the
state when in fact the \emph{state} is wrong, leads the agent to distrust the evidence it needs. This
rests on one seed and a manual inspection, and we never varied the NOTICE, so it is a hypothesis to
test. On wrong reads, by contrast, adding the contract helps Haiku (\RHaikuReadBothSuccess{} against
\RHaikuReadReadbackSuccess{}).

\paragraph{Practical implications.} These suggestions rest largely on the exploratory results. Of the
triggers here, read-back is the only one that reliably sees writes which echoed the right thing, at the
cost of extra reads and an expected-effect model, so it is worth adding where that effect can be
written down. Whether it pays depends on the model: it needs to act on true flags and to shrug off
false ones. On Qwen3-4B the combined check lowered success overall (\RExpQwenFourPThreeDiff{}), almost
all of it through false alarms on delayed writes, so where writes may become visible late, read-back
should wait or re-check before it flags. Keep recovery rather than abstaining when a wrong commit is
not much worse than an abandoned task; abstain where a wrong commit is costly. Flagging every step
hurt the small models on clean tasks and helped Haiku on silent faults, so its value depends on the
model and it is no substitute for a targeted trigger. Contracts are
cheap and catch implausible values, but they cannot see a write that echoed the right thing, and their
message should say what they found.
